
\section{Linguaggi e funzioni}
Esattamente come gli input sono codificati come stringhe su un alfabeto $\Sigma$, i problemi di decisione possono essere codificati come linguaggi, cioé sottoinsiemi di stringhe dell'alfabeto. \\
Una macchina di Turing risponde $"yes"$ a una istanza di un problema se la codifica dell'istanza é una stringa che appartiene al linguaggio. \\
Analogamente, i problemi che richiedono un output possono essere visti come funzioni. \\

\subsection{Linguaggi ricorsivi}
$L \subseteq \Sigma^*$ é \textbf{ricorsivo} sse \[\exists M :  M(x) =
\begin{cases}
	"yes" \qquad se \qquad x \in L \\
	"no" \qquad altrimenti
\end{cases}
\] \\
$M$ \textbf{decide} $L$. \\

\subsection{Linguaggi ricorsivamente enumerabili}
$L \subseteq \Sigma^*$ é \textbf{ricorsivamente enumerabile} sse \[\exists M :  M(x) =
\begin{cases}
	"yes" \qquad se \qquad x \in L \\
	\nearrow \qquad altrimenti
\end{cases}
\] \\
$M$ \textbf{accetta} $L$. \\

\subsection{Proposizione: i linguaggi ricorsivi sono ricorsivamente enumerabili}
\boxed{L \quad \textbf{ricorsivo} \implies L \quad \textbf{ricorsivamente enumerabile}} \\

\textbf{Prova:} \\
\[\exists M :  M(x) =
\begin{cases}
	"yes" \qquad se \qquad x \in L \\
	"no" \qquad altrimenti
\end{cases}
\] \\
Costruiamo $M' = (Q', \Sigma, \delta', s)$ aggiungendo un nuovo stato $q'$ che, nella nuova funzione di transizione $\delta'$, prende il posto dello stato terminale $"no"$. La funzione di transizione é uguale a quella di M in tutti gli altri casi. Il nuovo stato $q'$ sposta la testina a destra all'infinito. Formalmente:
\begin{itemize}
	\item $Q' = Q \cup \{q'\}$ \quad dove \quad $q' \notin Q$
	\item $\delta(q, \sigma) = ("no",  \tau, d) \implies \delta'(q, \sigma) = (q', \tau, d)  \qquad \forall q \in Q;  \forall \sigma, \tau \in \Gamma;  \forall d \in D$
	\item $\delta(q, \sigma) = (p,  \tau, d) \implies \delta'(q, \sigma) = (p, \tau, d)  \quad \forall q, p \in Q \setminus \{"no"\};  \forall \sigma, \tau \in \Gamma;  \forall d \in D$
	\item $\delta'(q', \sigma) = (q', \sigma, \rightarrow) \qquad \forall \sigma \in \Sigma \cup \{\sqcup\}$
\end{itemize}
\[M'(x) =
\begin{cases}
	"yes" \qquad se \qquad x \in L \\
	\nearrow \qquad altrimenti
\end{cases}
\] \\
\textbf{QED}

\subsection{Funzioni ricorsive}
La funzione totale $f : \Sigma^* \rightarrow {\Sigma_\sqcup}^*$ é \textbf{ricorsiva} sse \[\exists M :  \forall x \in \Sigma^* : M(x) = f(x) \] \\
$M$ \textbf{calcola} $f$. \\

\subsection{Funzioni ricorsivamente enumerabili}
La funzione parziale\footnote{Ricordando che le funzioni parziali non sono definite su tutto il dominio, e dunque esistono stringhe $x \in \Sigma^*$ per cui $f$ non é definita.}  $f : \Sigma^* \rightarrow {\Sigma_\sqcup}^*$ é \textbf{ricorsivamente enumerabile} sse \[\exists M :  \forall x \in \Sigma^* : \begin{cases}
	M(x) = y \qquad se \quad \exists y \in {\Sigma_\sqcup}^* : f(x) = y \\
	M(x) = \nearrow \qquad altrimenti
\end{cases}
\] \\

\subsection{Proposizione: una funzione é ricorsiva se e solo se é totale e ricorsivamente enumerabile}
\boxed{f \text{ é ricorsiva} \iff (f \text{ é totale} \land f \text{ é ricorsivamente enumerabile} )} \\

\textbf{Prova:} \\
$\implies$: Per definizione, $f$ é totale e	$\exists M : \forall x \in \Sigma^* \exists y \in {\Sigma_\sqcup}^* : M(x) = f(x) = y$ \\
$\impliedby$:  Siccome $f$ é totale,  $ \forall x \in \Sigma^* \exists y \in {\Sigma_\sqcup}^* : f(x) = y$.
Quindi il caso $M(x) = \nearrow$ non si applica mai. \\
\textbf{QED} \\


\subsection{Tesi di Church}
Le funzioni calcolabili corrispondono a quelle ricorsive, cioé calcolabili da una macchina di Turing. \\
Ogni tentativo ragionevole di formalizzare il concetto di funzione calcolabile (dove per ragionevole intendiamo numero di istruzioni finito e strutture dati finite) porta a una definizione equivalente a quella di macchina di Turing, con una perdita di efficienza al piú polinomiale. \\

Non é dimostrata, ma sono dimostrabili risultati particolari, come ad esempio:
\subsubsection*{Teorema: Equipotenza tra Macchine di Turing e RAM}
Se una RAM calcola una funzione $\phi$ in tempo $f(n)$, allora esiste una MdT $M$ a 7 nastri che calcola $\phi$ in tempo $O(f(n)^3)$. \\